\documentclass[../main.tex]{subfiles}

\begin{document}

\subsection{Python and NumPy Primer}
\label{sub:numpy-primer}

\textBox{%
This section summarizes the NumPy array conventions, linear-algebra operations, and indexing idioms used throughout the computational examples in these notes.
}

\subsubsection{Arrays and Basic Operations}

\addDef{NumPy Arrays}{%
The \texttt{ndarray} is the fundamental structure.
\begin{itemize}
  \item \textbf{Creation:} \texttt{np.array([1, 2])}, \texttt{np.zeros((m, n))}, \texttt{np.random.randn(m, n)}.
  \item \textbf{Indexing:} 0-indexed. \texttt{A[i, j]} for entries, \texttt{A[i, :]} for rows, and \texttt{A[:, j]} for columns.
\end{itemize}
}
\label{def:numpy-arrays}

\addRemark{%
\textbf{(Shapes and rank-1 arrays)}\enspace
\texttt{A.shape} returns \texttt{(m, n)}. A 1D array has shape \texttt{(n,)}, which is \textbf{not} equivalent to a column vector \texttt{(n, 1)}. These ``rank-1 arrays'' are a common source of broadcasting bugs. Use \texttt{x[:, None]} or \texttt{x.reshape(-1, 1)} to convert to 2D.
}

\addExcr{%
\exCompute
\begin{enumerate}
  \item Create $4 \times 4$ identity matrix. Verify \texttt{shape} and \texttt{dtype}.
  \item Extract row 2 and column 3 from a random $3 \times 3$ $\bA$.
  \item Compare \texttt{A @ x} for \texttt{x.shape == (3,)} and \texttt{x.shape == (3,1)}. Record the output shape in each case, then explain why the two representations behave differently under broadcasting.
\end{enumerate}
}

\subsubsection{Linear Algebra Operations}

\textBox{%
The core routines in \texttt{numpy.linalg} provide access to standard BLAS and LAPACK implementations:
}

\begin{center}
\begin{tabular}{ll}
\textbf{Function} & \textbf{Computes} \\ \hline
\texttt{A @ B} & Matrix product $\bA\bB$ \\
\texttt{np.linalg.solve(A, b)} & Solve $\bA\bx = \bb$ ($O(n^3)$) \\
\texttt{np.linalg.inv(A)} & Matrix inverse (prefer \texttt{solve} for linear systems) \\
\texttt{np.linalg.norm(x)} & Euclidean norm $\|\bx\|_2$ \\
\texttt{np.linalg.eig(A)} & Eigenvalues and eigenvectors \\
\texttt{np.linalg.svd(A)} & Singular Value Decomposition \\
\texttt{np.linalg.qr(A)} & QR Factorization \\
\texttt{np.linalg.cond(A)} & Condition number $\kappa_2(\bA)$ \\
\end{tabular}
\end{center}

\addRemark{%
\textbf{(Vectorization and performance)}\enspace
Explicit \texttt{for} loops in Python are slow. Push heavy lifting into NumPy/BLAS routines by \textbf{vectorizing} operations. For example, use \texttt{A @ B} instead of triple-nested loops.
}

\addRemark{%
\textbf{(Linear solve versus inversion)}\enspace
\texttt{np.linalg.solve(A, b)} is faster and more stable than \texttt{np.linalg.inv(A) @ b}. Explicit inversion is almost never desirable in scientific computation.
}

\addRemark{%
\textbf{(Computed equality)}\enspace Use \texttt{np.allclose(x, y)} or \texttt{np.isclose(x, y)} for computed floating-point quantities. Exact equality is appropriate for integers, booleans, and deliberately exact structural checks.
}

\addRemark{%
\textbf{(Sparse matrices)}\enspace Large discretized PDE and graph problems should usually use \texttt{scipy.sparse}, not dense NumPy arrays. Use \texttt{scipy.sparse.linalg} for sparse solves, eigenvalue computations, CG, and GMRES.
}

\addExcr{%
\exCompute
\begin{enumerate}
  \item Solve a $2 \times 2$ system and verify the residual $\|\bA\bx - \bb\|$.
  \item Profile \texttt{solve} vs \texttt{inv @ b} for $n=1000$.
  \item Compare \texttt{np.linalg.eig} and \texttt{np.linalg.eigh} (optimized for symmetric matrices).
\end{enumerate}
}

\subsubsection{Key Idioms}

\begin{itemize}
  \item \textbf{Element-wise:} \texttt{*} and \texttt{**} are entry-wise. Matrix powers use \texttt{np.linalg.matrix\_power}.
  \item \textbf{Broadcasting:} Automatic dimension expansion. \texttt{A * v} scales columns of $\bA$ by entries of $v$.
  \item \textbf{Aliases, Views, Copies:} \texttt{B = A} makes a second name for the same array. Slices are usually views. Use \texttt{B = A.copy()} for independent data.
  \item \textbf{Sparse construction:} use \texttt{scipy.sparse.diags} for banded matrices and \texttt{scipy.sparse.csr\_matrix} for efficient arithmetic.
\end{itemize}

\addExcr{%
\exCompute
\begin{enumerate}
  \item Let $\bA = \begin{pmatrix} 1 & 2 \\ 3 & 4 \end{pmatrix}$. Compare \texttt{A * A} and \texttt{A @ A}.
  \item Demonstrate the ``view pitfall'' by modifying a slice of a matrix.
  \item Use broadcasting to perform diagonal scaling $\bA\bD$ without constructing $\bD = \text{diag}(d)$.
\end{enumerate}
}

\addReading{%
Heath gives a broad introduction to scientific computing at roughly the level of these notes \citep{heath2018}. For the effects of floating-point arithmetic, Higham is the standard detailed reference \citep{higham2002}. Meyer provides an accessible bridge from undergraduate linear algebra to matrix computation \citep{meyer2000}.
}

\end{document}
